\documentclass[10pt,a4paper]{amsart}
\usepackage[margin=19mm]{geometry}
\usepackage{amsmath,amssymb,mathtools}
\usepackage[scr=rsfs,cal=euler]{mathalfa}
\usepackage{xfrac}
\usepackage[unicode,hidelinks,bookmarks=false]{hyperref}
\newcommand{\ie}{i.e.\ }
\newcommand{\cf}{cf.\ }
\newcommand{\resp}{respectively\ }
\newcommand{\Subsection}{\subsection}
\providecommand{\nobreakdash}{\nobreak\hbox{-}}
\setlength{\emergencystretch}{2em}
\multlinegap0pt
\usepackage{amssymb}
\usepackage{xcolor}
\newtheorem{theorem}{Theorem}
\newtheorem{corollary}{Corollary}
\newtheorem{lemma}{Lemma}
\title{Umbral calculus over a vector space}
\author{Abdullah Alharthi, Eugene Lytvynov}
\date{Source: arXiv:2609.10030v1 (CC BY 4.0)}
\begin{document}
\maketitle

\noindent\emph{Source and licence.} Umbral calculus over a vector space. Version arXiv:2609.10030v1 (CC BY 4.0), distributed by arXiv under the Creative Commons Attribution 4.0 International licence, http://creativecommons.org/licenses/by/4.0/, as recorded in the arXiv licence metadata for this exact version and shown on the abstract page https://arxiv.org/abs/2609.10030v1. Full licence notice: \texttt{licenses/arxiv-2609.10030-CC-BY-4.0.txt}.\par\smallskip

\noindent\textit{Editorial scope.} Counted unit: the article's theorem dzerste7kui7 (source0.tex lines 1131--1136). The article's complete original proof of it is delivered with it (source0.tex lines 1172--1315, one \texttt{proof} environment of 144 lines, which the article places away from the statement). No mathematics is written, repaired or re-derived: the statement and the proof are the article's own lines, in the article's own notation, and no retained line is substituted or re-flowed.  Retained uncounted as the source-local context the proof needs: lines 461--463; lines 504--512; lines 506--507; lines 568--577; lines 606--619; lines 649--652; lines 675--677; lines 769--771; lines 1060--1062.  Nothing was omitted from any retained span, so the package carries no omission record.  No retained line contains a citation, so the wrapper carries no bibliography and no bibliography entry.  No retained line contains a figure environment, an image-inclusion command or drawing code, so no image is delivered and the package declares no asset dependency.  Editorial formatting only, in addition: an AMS \texttt{amsart} preamble that restates the article's own declarations and macros used by the retained lines, namely the environments \texttt{theorem}, \texttt{corollary}, \texttt{lemma} on separate counters, together with the standard packages those lines need.  The retained environments print in this document as Theorem 1, Corollary 1, Lemma 1 in order of appearance, whereas the article prints the same items with the numbering of its own full document.  The complete native source of the article as distributed in the arXiv e-print for this exact version is bundled unchanged as \texttt{sources/209869/source0.tex}.
\begin{equation}\label{ge4ew34wq43w}
\big(Q^{k}\langle P_n(\cdot),f_n\rangle\big)(\omega)=(n)_k\langle P_{n-k}(\omega),f_n\rangle,
\end{equation}

\begin{theorem}[Operator expansion]\label{bgnfm758} Let $P\in\mathbb P(V)$ and let $(P_n(\omega))_{n=0}^\infty$ be the corresponding polynomial sequence. Assume that 
  $P_0=1$ and $P_n (0)=0$ for all $n\geq 1$.    Assume that the corresponding lowering operators $Q(\theta)$  ($\theta\in V^*$) are shift-invariant. Let $T \in \mathcal L(\mathcal P  (V^*))$. Then $T$ is shift-invariant, i.e., $T\in\mathcal{SI}(\mathcal P  (V^*))$, if and only if there exist $T_0\in \mathbb F$ and $T_k\in (V^{\odot  k})^*$ for $k\in \mathbb N$ such that
\begin{equation}\label{vgxds5ese5}T=T_0\mathbf 1+\sum_{k=1}^\infty \langle T_k,
Q^{k}\rangle .\end{equation}
In the latter case,  
\begin{equation}\label{gxes5seasew}T_0=T 1,\quad \langle 
T_k, f_k \rangle =\frac1{k!}\big(T\langle P_k(\cdot),f_k\rangle\big)(0)\quad\text{for $k\in\mathbb N$ and $f_k\in V^{\odot   k}$}.
\end{equation} 
\end{theorem}

\begin{equation}T=T_0\mathbf 1+\sum_{k=1}^\infty \langle T_k,
Q^{k}\rangle .\end{equation}

\begin{corollary}\label{567ty7} 
  

{\rm (i)\/} For  $T \in \mathcal {SI}(\mathcal P  (V^*))$ of the form \eqref{vgxds5ese5}, define a formal tensor power series $(\mathcal IT)(v)\in \mathcal F(V,\mathbb F)$ by
$(\mathcal IT)(v):=T_0+\sum_{k=1}^\infty \langle T_k,v^{\otimes k}\rangle$.
Then $\mathcal I:\mathcal{SI}(\mathcal P  (V^*))\to \mathcal F(V,\mathbb F)$
 is an algebra isomorphism. In particular, any two shift-invariant operators commute. 

{\rm (ii)\/} A shift-invariant operator $T\in \mathcal{SI}(\mathcal P    (V^*))$  is invertible if and only if $T1\ne0$. In the latter case, $T^{-1}$ is also shift-invariant. 
\end{corollary}

\begin{corollary}\label{xesa4aq46q327y} 
{\rm (i)} Assume that the  lowering operators $Q(\theta)$  ($\theta\in V^*$) are as in Theorem~\ref{bgnfm758}.  Let $G \in \mathcal L\big(\mathcal P  (V^*),\mathcal P  (V^*,V)\big)$. Then  $G\in\mathcal{SI}\big(\mathcal P  (V^*),\mathcal P(V^*,V)\big)$ if and only if there exist $G_0\in V$ and $G_k\in \mathcal L(V^{\odot k},V)$ for $k\in \mathbb N$ such that
\begin{equation}\label{utf6d6eid}
G=G_0\mathbf 1+\sum_{k=1}^\infty G_k
Q^{k}.\end{equation}
In the latter case,  
$G_0=G 1$ and $ G_k f_k  =\frac1{k!}\big(G\langle P_k(\cdot),f_k\rangle\big)(0)$ for $k\in\mathbb N$ and $f_k\in V^{\odot   k}$.

{\rm (ii)} For   $G\in\mathcal{SI}\big(\mathcal P  (V^*),\mathcal P(V^*,V)\big)$ of the form \eqref{utf6d6eid}, define a formal tensor power series $(\mathcal JG)(v)\in \mathcal F(V,V)$ by
$(\mathcal JG)(v)=G_0+\sum_{k=1}^\infty G_kv^{\otimes k}$.
Then the map
$$\mathcal J:\mathcal{SI}\big(\mathcal P  (V^*),\mathcal P(V^*,V)\big)\to \mathcal F(V,V)$$
 is bijective.
\end{corollary}

\begin{lemma}
\label{gfyrdr6ew64ui}
Let $(P_n(\omega))_{n=0}^\infty$ be a binomial sequence. Then $P_n(0)=0$ for all $n\in\mathbb N$.
\end{lemma}

\begin{equation}\label{trse5sw5w}
\sum_{n=0}^\infty \frac{1}{n!}\,\langle P_n(\omega),v^{\otimes n}\rangle= \exp\big[\langle \omega,B(v)\rangle\big],\quad \omega \in V^*,
\end{equation}

\begin{equation}\label{ctxsts}
\big(\mathcal I\langle T_k,Q^{k}\rangle \big)(v)=\langle T_k,L(v)^{\otimes k}\rangle.
 \end{equation}

\begin{equation}
\langle S_n(\omega),v_1\odot\dots\odot v_n\rangle=\sum_{\pi\in\mathcal P(n)}\prod_{\mathcal A\in\pi}\Big(\big\langle\omega,\tilde B_{|\mathcal A|}\underset{i\in \mathcal A}\odot v_i\big\rangle+\big\langle \tilde\alpha_{|\mathcal A|},\underset{i\in \mathcal A}\odot v_i\big\rangle\Big).\label{cxrtstset55}
\end{equation}

\begin{theorem}\label{dzerste7kui7} Under the above assumptions, we have, for each $\zeta\in V$, $p\in\mathcal P(V^*)$ and $\omega\in V^*$,
\begin{equation}\label{dtrd6t5swe5u4wu5}
\big(R(\zeta)p\big)(\omega)=
\big\langle\omega, (\Phi_\zeta(D)p)(\omega)\big\rangle+ (\Gamma_\zeta(D)p)(\omega). 
\end{equation}
\end{theorem}

\begin{proof}[Proof of Theorem~\ref{dzerste7kui7}] We divide the proof into several 
steps.

{\it Step 1}. For each $k\in\mathbb N$, let us fix  $T_k\in(V^{\odot k})^*$. For $k,n\in\mathbb N$ with $n\ge k$ and any $v_1,\dots,v_n\in V$, formula \eqref{ge4ew34wq43w} with $Q=D$ implies
$$ \frac1{k!}\,\big(\langle T_k ,D^{k}\rangle\langle\cdot ^{\otimes n},v_1\odot\dots\odot v_n\rangle\big)(\omega)=\sum_{{\mathcal A\subset\{1,\dots,n\}}\atop{|\mathcal A|=k}}\big\langle T_k ,\underset{i\in \mathcal A}\odot v_i\big\rangle\big\langle\omega^{\otimes(n-k)},\underset{j\in \mathcal A^c}\odot v_j\big\rangle,$$
where $\mathcal A^c=\{1,\dots,n\}\setminus\mathcal A$. Hence, for any constants $c_1,\dots,c_n\in\mathbb F$, we obtain
$$ \frac1{k!}\,\bigg(\langle T_k, D^{k}\rangle\prod_{i=1}^n(\langle\cdot,v_i\rangle+c_i)\bigg)(\omega)=\sum_{{\mathcal A\subset\{1,\dots,n\}}\atop{|\mathcal A|=k}}\big\langle T_k, \underset{i\in \mathcal A}\odot v_i\big\rangle\prod_{j\in\mathcal A^c}(\langle\omega,v_j\rangle+c_j).$$
Therefore,
\begin{equation}\label{dxters5ewu}
\bigg(\sum_{k=1}^\infty \frac 1{k!}\,\langle T_k, D^{k}\rangle\prod_{i=1}^n(\langle\cdot,v_i\rangle+c_i)\bigg)(\omega) =\sum_{{\mathcal A\subset\{1,\dots,n\}}\atop{\mathcal A\ne\varnothing}}\big\langle T_{|\mathcal A|}, \underset{i\in \mathcal A}\odot v_i\big\rangle\prod_{j\in\mathcal A^c}(\langle\omega,v_j\rangle+c_j).\end{equation}

{\it Step 2}. Let $\pi=\{\mathcal A_1,\dots,\mathcal A_k\}\in\mathcal P(n)$. We write $\rho\subset\pi$, $\rho\ne\varnothing$ if $\rho=\{\mathcal A_{i_1},\dots,\mathcal A_{i_m}\}$, where $\{i_1,\dots,i_m\}$ is a nonempty subset of $\{1,\dots,k\}$. By \eqref{dxters5ewu}, we have for each $\pi\in\mathcal P(n)$,
\begin{align}
&\bigg(\sum_{k=1}^\infty \frac 1{k!}\,\langle T_k, D^{k}\rangle \prod_{\mathcal A\in\pi}\Big(
\big\langle \cdot,\tilde B_{|\mathcal A|}\underset{i\in \mathcal A}\odot v_i\big\rangle+\big\langle \tilde\alpha_{|\mathcal A|},\underset{i\in \mathcal A}\odot v_i\big\rangle\Big)\bigg)(\omega)\notag\\
&\quad=\sum_{\rho\subset\pi\atop \rho\ne\varnothing}\Big\langle T_{|\rho|},
\underset{\mathcal A\in\rho}\odot\Big(\tilde B_{|\mathcal A|} \underset{i\in\mathcal A}\odot v_i\Big)\Big\rangle
\prod_{\mathcal B\in\pi\atop\mathcal B\not\in \rho}
\bigg(
\big\langle \omega,\tilde B_{|\mathcal B|}\underset{j\in \mathcal B}\odot v_j\big\rangle+\big\langle \tilde\alpha_{|\mathcal B|},\underset{j\in \mathcal B}\odot v_j\big\rangle\bigg).\label{xersa5y}
\end{align}
By \eqref{cxrtstset55} and \eqref{xersa5y}, we have, for  $v_1,\dots,v_n\in V$,
\begin{align}
&\bigg(\sum_{k=1}^\infty \frac 1{k!}\,\langle T_k, D^{k}\rangle\langle S_n(\cdot),v_1\odot\dots\odot v_n\rangle\bigg)(\omega)\notag\\ 
&\quad=\sum_{\pi\in\mathcal P(n)}\sum_{\mathcal A\in\pi}
\bigg(\sum_{\rho\in\mathcal P(\mathcal A)}\Big\langle T_{|\rho|},\underset{\mathcal B\in\rho}{\odot}\big(\tilde B_{|\mathcal B|}\underset{i\in\mathcal B}\odot v_i\big)\Big\rangle \bigg)
\prod_{{\mathcal C\in\pi}\atop{\mathcal C\ne\mathcal A}}\Big(\langle\omega,\tilde B_{|\mathcal C|}\underset{j\in\mathcal C}\odot v_j \rangle+\langle\tilde\alpha_{|\mathcal C|},
\underset{j\in\mathcal C}\odot v_j\rangle\Big).\label{vytdyd}
\end{align}
(In  formula \eqref{vytdyd}, if $\pi$ contains a single set $\mathcal A=\{1,\dots,n\}$, the product over $\mathcal C\in\pi$ with $\mathcal C\ne\mathcal A$ is supposed to be equal to 1.)

{\it Step 3}. Denote by $(P_n(\omega))_{n=0}^\infty$ the binomial sequence corresponding to the Sheffer sequence $(S_n(\omega))_{n=0}^\infty$, i.e., $(P_n(\omega))_{n=0}^\infty$ has generating function \eqref{trse5sw5w}. Then, by formula~\eqref{vytdyd}, 
\begin{align}
&\bigg(\sum_{k=1}^\infty \frac 1{k!}\,\langle T_k, D^{k}\rangle\langle P_n(\cdot),v_1\odot\dots\odot v_n\rangle\bigg)(\omega)\notag\\ 
&\quad=\sum_{\pi\in\mathcal P(n)}\sum_{\mathcal A\in\pi}
\bigg(\sum_{\rho\in\mathcal P(\mathcal A)}\Big\langle T_{|\rho|},\underset{\mathcal B\in\rho}{\odot}\big(\tilde B_{|\mathcal B|}\underset{i\in\mathcal B}\odot v_i\big)\Big\rangle \bigg)
\prod_{{\mathcal C\in\pi}\atop{\mathcal C\ne\mathcal A}}\langle\omega,\tilde B_{|\mathcal C|}\underset{j\in\mathcal C}\odot v_j \rangle.\notag
\end{align}
Therefore,
\begin{equation}\label{fdre654w6}
\bigg(\sum_{k=1}^\infty \frac 1{k!}\,\langle T_k, D^{k}\rangle\langle P_n(\cdot),v_1\odot\dots\odot v_n\rangle\bigg)(0)= \sum_{\rho\in\mathcal P(n)}\Big\langle T_{|\rho|},\underset{\mathcal B\in\rho}{\odot}\big(\tilde B_{|\mathcal B|}\underset{i\in\mathcal B}\odot v_i\big)\Big\rangle .
\end{equation}



{\it Step 4}. 
For each $\zeta\in V$ and $n\in\mathbb N$, we define $\widehat B_{n}(\zeta)\in\mathcal L(V^{\odot n},V)$ and $\widehat\alpha_{n}(\zeta)\in (V^{\odot n})^*$ by 
\begin{equation}\label{dy6de6ir5}
\widehat B_{n}(\zeta)f_n:=\tilde B_{n+1}(\zeta\odot f_n),\quad \langle\widehat\alpha_{n}(\lambda),f_n\rangle:=\langle\tilde\alpha_{n+1},\zeta\odot f_n\rangle,\quad f_n\in V^{\odot n}.
\end{equation}
Then, by formula \eqref{cxrtstset55}, we have, for $\zeta,v\in V$ and $n\in\mathbb N$,
\begin{align}
&\langle S_{n+1}(\omega),\zeta\odot v^{\otimes n}\rangle=\big(\langle\omega, B_1\zeta\rangle+\langle \alpha_1,\zeta\rangle\big)\langle S_n(\omega),v^{\otimes n}\rangle\notag\\
&\quad+\sum _{\pi\in\mathcal P(n)}\sum_{\mathcal A\in\pi} \big(\langle\omega,\widehat B_{|\mathcal A|}(\zeta)  v^{\otimes|\mathcal A|}\rangle+\langle \widehat\alpha_{|\mathcal A|}(\zeta),v^{\otimes|\mathcal A|}\rangle\big) 
\prod_{{\mathcal B\in\pi}\atop{\mathcal B\ne\mathcal A}}\big(\langle\omega,\tilde B_{|\mathcal B|}v^{\otimes|\mathcal B|}\rangle+\langle\tilde\alpha_{|\mathcal B|},v^{\otimes|\mathcal B|}\rangle\big).\label{cdtre6u}
\end{align}
Note also that
\begin{equation}
\langle S_{1}(\omega),\zeta\rangle= \langle\omega, B_1\zeta\rangle+\langle \alpha_1,\zeta\rangle.\label{vdtrs6ue}
\end{equation}

 {\it Step 5}. Fix $\zeta\in V$. Assume that $\Theta_k(\zeta)\in \mathcal L(V^{\odot k},V)$ and $\Xi_k(\zeta)\in(V^{\odot k})^*$
($k\in\mathbb N$) are 
 such that, for each  $v\in V$ and $n\in\mathbb N$,
\begin{align}
&\sum_{\rho\in\mathcal P(n)}  \Theta_{|\rho|}(\zeta)\bigg(\underset{\mathcal B\in\rho}{\odot}\big(\tilde B_{|\mathcal B|}v^{\otimes |\mathcal B|}\big)\bigg)=\widehat B_{n}(\zeta)v^{\otimes n} ,\label{vfyxreasry5ay45}\\
&\sum_{\rho\in\mathcal P(n)} \Big\langle \Xi_{|\rho|}(\zeta), \underset{\mathcal B\in\rho}{\odot}\big(\tilde B_{|\mathcal B|}v^{\otimes |\mathcal B|}\big)\Big\rangle =\langle \widehat\alpha_{n}(\zeta),v^{\otimes n}\big\rangle.
\label{bvyuf7}\end{align}
Then, by \eqref{vytdyd}, \eqref{vfyxreasry5ay45} and \eqref{bvyuf7}, we obtain, for $\omega\in V^*$ and $v\in V$,
\begin{align}
&\bigg(\sum_{k=1}^\infty \frac 1{k!}\,\langle \omega \Theta_k(\zeta)+\Xi_k(\zeta), D^{k}\rangle
\langle S_n(\cdot),v^{\otimes n}\rangle\bigg)(\omega)\notag\\ 
&\quad=\sum_{\pi\in\mathcal P(n)}\sum_{\mathcal A\in\pi}
\big(\langle\omega,\widehat B_{|\mathcal A|}v^{\otimes|\mathcal A|}\rangle+\langle\widehat \alpha_{|\mathcal A|},v^{\otimes |\mathcal A|}\rangle \big)
\prod_{{\mathcal B\in\pi}\atop{\mathcal B\ne\mathcal A}}\Big(\langle\omega,\tilde B_{|\mathcal B|}v^{\otimes|\mathcal B|} \rangle+\langle\tilde\alpha_{|\mathcal B|},
v^{\otimes|\mathcal B|}\rangle\Big).\label{vcfxtsxte5su5}
\end{align}

By \eqref{cdtre6u} and \eqref{vcfxtsxte5su5}, we have, for each $p\in\mathcal P(V^*)$, $\omega\in V^*$ and $v\in V$,
\begin{align}
(R(\zeta)p)(\omega)&=\big(\langle\omega, B_1\zeta\rangle+\langle \alpha_1,\zeta\rangle\big)p(\omega)+\bigg\langle\omega, \bigg(\sum_{k=1}^\infty \frac 1{k!}\,\Theta_k(\zeta)D^kp\bigg)(\omega)\bigg\rangle\notag\\
&\quad+\bigg(\sum_{k=1}^\infty \frac 1{k!}\,\langle \Xi_k(\zeta), D^k\rangle p\bigg)(\omega).\label{vgdctrs5ea43y7}
\end{align}


{\it Step 6}. In view of \eqref{fdre654w6}, conditions \eqref{vfyxreasry5ay45},  \eqref{bvyuf7}   can be reformulated as follows: for each 
$\zeta\in V$, $v\in V$ and $n\in\mathbb N$, 
\begin{align}
&\bigg(\sum_{k=1}^\infty \frac 1{k!}\,\Theta_k(\zeta) D^{k} \langle P_n(\cdot),v^{\otimes n}\rangle\bigg)(0)=
\widehat B_{n}(\zeta)v^{\otimes n},\label{vdrste5s5sw5}\\
&\bigg(\sum_{k=1}^\infty \frac 1{k!}\,\langle \Xi_k(\zeta), D^{k} \rangle\langle P_n(\cdot),v^{\otimes n}\rangle\bigg)(0)=\langle \widehat\alpha_{n}(\zeta),v^{\otimes n}\rangle.
\label{dzrwaw}\end{align}


By \eqref{ge4ew34wq43w} and Lemma \ref{gfyrdr6ew64ui}, we have, for any $T_k\in(V^{\odot k})^*$ ($k\in\mathbb N$), $v\in V$ and $n\in\mathbb N$,
\begin{equation*}
\bigg(\sum_{k=1}^\infty \frac 1{k!}\,\langle T_k, Q^{k}\rangle\langle P_n(\cdot),v^{\otimes n}\rangle\bigg)(0)=\langle T_n,v^{\odot n}\rangle.
\end{equation*}
Therefore,
\begin{align*}
\bigg(\sum_{k=1}^\infty \frac 1{k!}\,\widehat B_{k}(\zeta)Q^{k}\langle P_n(\cdot),v^{\otimes n}\rangle\bigg)(0)&=\widehat B_{n}(\zeta)v^{\odot n},\\
\bigg(\sum_{k=1}^\infty \frac 1{k!}\,\langle  \widehat\alpha_{k}(\zeta), Q^{k}\rangle\langle P_n(\cdot),v^{\otimes n}\rangle\bigg)(0)&=\langle  \widehat\alpha_{n}(\zeta),v^{\odot n}\rangle
\end{align*}
Hence, conditions \eqref{vdrste5s5sw5}, \eqref{dzrwaw} are satisfied provided the following equalities hold true:
\begin{align}
\sum_{k=1}^\infty \frac 1{k!}\,  \Theta_k(\zeta) D^{k} &=\sum_{k=1}^\infty \frac 1{k!}\,\widehat B_{k}(\zeta) Q^{k}\label{cdaewa4yw6u7j8},\\
\sum_{k=1}^\infty \frac 1{k!}\,  \langle \Xi_k(\zeta), D^{k}\rangle &=\sum_{k=1}^\infty \frac 1{k!}\,\langle \widehat \alpha_{k}(\zeta), Q^{k}\rangle.
\label{vyjdk7r8}\end{align}

In view of Corollaries~\ref{567ty7} (i) and \ref{xesa4aq46q327y} (ii) with $Q=D$, conditions \eqref{cdaewa4yw6u7j8}, \eqref{vyjdk7r8} can be written as follows:
\begin{align}
\sum_{k=1}^\infty \frac 1{k!}\,  \Theta_k(\zeta) v^{\otimes k}&=\sum_{k=1}^\infty \frac 1{k!}\,\mathcal J\big(\widehat B_{k}(\zeta) Q^{k}\big)(v),\label{vcyd6i5}\\
\sum_{k=1}^\infty \frac 1{k!}\, \langle \Xi_k(\zeta), v^{\otimes k}\rangle&=\sum_{k=1}^\infty \frac 1{k!}\,\mathcal I\big(\langle \widehat \alpha_{k}(\zeta), Q^{k}\rangle\big)(v).\label{dsthe}
\end{align}

We easily conclude from \eqref{ctxsts} that 
\begin{align}
&\sum_{k=1}^\infty \frac 1{k!}\,\mathcal J\big(\widehat B_{k}(\zeta) Q^{k}\big)(v)=
\sum_{k=1}^\infty \frac 1{k!}\,\widehat B_{k}(\zeta) L(v)^{\otimes k}
=\sum_{k=1}^\infty \frac 1{k!}\,\tilde B_{k+1}(\zeta\odot L(v)^{\otimes k})\notag\\
&\quad= \sum_{k=2}^\infty\frac{k}{k!}\,\tilde B_k(\zeta\odot L(v)^{\otimes(k-1)})=\breve B'(v;\zeta)\circ L(v). \label{vftr6e645}
\end{align}
Here, we used the formal tensor power series 
$$\breve B(v):=\sum_{k=2}^\infty \frac1{k!}\tilde B_kv^{\otimes k}=B(v)-B_1v.$$ 
Similarly,
\begin{equation}
\sum_{k=1}^\infty \frac 1{k!}\,\mathcal I\big(\widehat \alpha_{k}(\zeta) Q^{k}\big)(v)=\breve \alpha'(v;\zeta)\circ L(v),
\label{ryd6ue5i7}
\end{equation}
where 
$$\breve \alpha (v):=\sum_{k=2}^\infty \frac1{k!}\langle \tilde \alpha _k,v^{\otimes k}\rangle=\alpha(v)-\langle \alpha_1,v\rangle.$$ 

Thus, conditions \eqref{vcyd6i5}, \eqref{dsthe} are satisfied if we choose $\Theta_k(\zeta)$ and $\Xi_k(\zeta)$ as follows:
\begin{equation}\label{vcyd6uei6}
\sum_{k=1}^\infty \frac 1{k!}\,  \Theta_k(\zeta) v^{\otimes k}=\breve B'(v;\zeta)\circ L(v),\quad \sum_{k=1}^\infty \frac 1{k!}\,  \langle \Xi_k(\zeta), v^{\otimes k}\rangle=\breve\alpha'(v;\zeta)\circ L(v).\end{equation}

 Denote $\breve\Phi_\zeta(v):=\breve B'(v;\zeta)\circ L(v)$ and $\breve\Gamma_\zeta(v):=\breve \alpha'(v;\zeta)\circ L(v)$. Then, by \eqref{vgdctrs5ea43y7} and \eqref{vcyd6uei6}, we obtain, for 
 $p\in\mathcal P(V^*)$, $\zeta\in V$ and $\omega\in V^*$,
\begin{align*}
(R(\zeta)p)(\omega)&=\big(\langle\omega, B_1\zeta\rangle
+\langle\alpha_1,\zeta\rangle\big)p(\omega)+\langle\omega,(\breve\Phi_\zeta(D)p)(\omega)\rangle+(\breve\Gamma_\zeta(D)p)(\omega)\\
&=\langle\omega,(\Phi_\zeta(D)p)(\omega)\rangle+(\Gamma_\zeta(D)p)(\omega).\qedhere
\end{align*}
\end{proof}

\end{document}
